% 第 6 章 已移动状态（由 tools/convert.py 从原 content/ 自动生成，勿手改）

虽然生成的或实现的移动语义通常不错，但至少应该了解可能的情况:移动操作会将对象置为C++标准库不支持的状态，或者破坏不变量。\par

这章中，我们根据C++标准库，来了解“无效”状态的定义，即已移动的对象所处的状态。\par


\section{已移动对象的要求和状态}

移动之后，已移动的对象既没有部分销毁，也没有完全销毁。析构函数还未调用，直到已移动对象的生命周期结束时才调用。因此，析构函数至少要没问题。\par

然而，C++标准库为可移动类型提供了更多的保证。已移动对象处于“有效但未定义的状态”，所以可以像使用任何不知道其值的类型的对象一样使用已移动对象。像使用该类型的非\textit{const}引用形参，而不知道所传递对象的值一样。例如，可以做的不仅仅是销毁已移动对象，还可以使用移动语义来实现排序和可变序列算法。\par

为了更详细地理解如何处理已移动对象，最好区分与之相关的信息:\par

\begin{itemize}
  \item C++标准库安全地使用已移动对象有什么要求?
  \item 给已移动对象什么样的保证，才能让这些类型的用户知道如何使用?
\end{itemize}

至少应该满足C++标准库的要求。\par

\mySmallsection{6.1.1 已移动对象的状态要求}

C++标准库对已移动对象的要求没有什么特别的。对于任何函数，为传递的类型和对象制定的需求，也适用于任何内部传递或是移动来的对象。\par

已移动对象必须能够销毁。此外，必须能够将新值赋给已移动对象。\par

例如，考虑如何交换两个对象\textit{a}和\textit{b}(这可能是排序操作的一部分)。交换的实现通常如下(见图6.1):\par

\begin{itemize}
  \item 将\textit{a}移动到一个新的临时对象\textit{tmp}(这样\textit{a}就变成了已移动的对象)。
  \item 将\textit{b}移动赋值给\textit{a}(这样\textit{b}就成为了已移动对象)。
  \item 将\textit{tmp}移动赋值给\textit{b}(这样\textit{tmp}就成了已移动对象)。
  \item 销毁已移动对象\textit{tmp}。
\end{itemize}

\myGraphic{0.9}{images/chapter6/1}{}

注意，可以通过向两个参数传递相同的对象来实现自交换。这种情况下，可以将状态为已移动的对象赋值给它自己。\par

所以，对于已移动对象，有同样的要求，适用于所有对象:\par

\begin{itemize}
  \item 必须能够销毁已移动对象。
  \item 必须能够为已移动对象赋值。
  \item 应该能够复制、移动或将已移动对象赋值给另一个对象。
\end{itemize}

已移动对象还能够处理特殊操作的额外需求。例如排序，所以必须支持<操作符，或有相应的排序标准，这也适用于已移动的对象。您可能会争辩说，排序算法中应该知道哪个对象移动了，这样就可以避免比较它，但C++标准库并不要求这样做。顺便一说，您可以传递已移动对象，对它们进行排序，只要支持所有必需的操作就可以了。\par

注意，C++标准库的定义需要的不仅仅是Alexander Stepanov和Paul McJones在《编程元素》一书的内容。\par

对于C++标准库中使用的所有对象和类型，应该确保已移动对象也支持所调用函数的要求。\par

\mySmallsection{6.1.2 保证已移动对象的状态}

已移动对象的保证定义在使用哪些代码。通常，类的设计者决定给出哪些保证，并且总在生命周期结束时销毁对象。对于移动状态，必须提定义良好的析构函数。\par

通常，会有更多要保障的。对于C++标准库的要求，支持基本操作，如复制、赋值相同类型的对象通常就足够了。然而，用户通常希望可以使用处理所有其他方式来为对象赋值。因此，保证可以将新值“赋值”给已移动对象就非常有用。\par

考虑给标准字符串\textit{s}赋新值的不同方法:\par

\begin{codebox}[cpp]
s = "hello"; // assign ”hello”
s.assign(s2); // assign the state of s2
s.clear(); // assign the empty state
std::cin >> s; // assign the next word from standard input
std::getline(myfile, s); // assign the next line from a file
\end{codebox}

例如，下面的循环是将流逐行读入vector对象的常用方法:\par

\begin{codebox}[cpp]
std::string row;
while (std::getline(myStream, row)) {
  coll.push_back(std::move(row)); // move the line into the vector
}
\end{codebox}

另一个例子中，可以在处理容器的泛型代码中实现以下语句:\par

\begin{codebox}[cpp]
foo(std::move(obj)); // pass obj to foo()
obj.clear(); // ensure the object is empty afterwards
\end{codebox}

类似的代码可以用于释放指针所使用的对象的内存:\par

\begin{codebox}[cpp]
draw(std::move(up)); // the unique pointer might or might not give up ownership
up.reset(); // ensure we give up ownership and release any resource
\end{codebox}

通过声明，已移动对象处于有效但未定义的状态，C++标准为其库中的操作进行了保证。\par

例如，以下是根据C++标准定义的行为:\par

\begin{codebox}[cpp]
std::stack<int> stk;
...
foo(std::move(stk)); // stk gets unspecified state
stk.push(42);
... // do something else without using stk
int i = stk.top();
assert(i == 42); // should never fail
\end{codebox}

虽然不知道通过\textit{std::move()}将\textit{stk}传递给\textit{foo()}后的值，但只要使用它来保存\textit{42}，直到再次需要，都可以将它当做有效的堆栈。\par

当然，所提供的保证能达到什么程度取决于您，但要确保您的类型的用户知道这些保证。通常，希望您的类型提供与C++库相同的保证，这样就可以像使用该类型的任何其他对象一样使用已移动对象了。\par

\mySmallsection{6.1.3 破坏不变量}

C++标准库定义了所有处于“有效但未定义状态”的移动对象的含义如下:\par

\textit{对象的值没有指定，除非满足对象的不变量，并且对象上的操作按照指定的类型进行}\par

不变量是适用于所有对象。有了这个保证，可以假定已移动的对象处于某种状态，这意味着它的不变量没有破坏。可以像使用非\textit{const}引用参数一样使用对象，而不知道传递的参数的任何信息:可以调用任何没有约束或前置条件的操作，并且效果/结果是为该类型的其他对象指定。\par

例如，对于已移动的字符串，可以执行没有先决条件的操作:\par

\begin{itemize}
  \item 查询大小
  \item 打印输出
  \item 遍历字符
  \item 转换为C字符串
  \item 分配新值
  \item 添加字符
\end{itemize}

此外，所有这些函数具有语义:\par

\begin{itemize}
  \item 返回的大小可以安全地调用索引操作符。
  \item 返回的大小与迭代时的字符数相匹配。
  \item 遍历字符时，打印的字符与字符序列匹配。
  \item 添加一个字符将把该字符放在值的末尾(不知道是否还有其他字符在前面)。
\end{itemize}

开发者可以利用这些保证(参见std::stack<>的例子)。\par

有时候，必须确保不破坏不变量，否则默认的特殊移动成员函数无法正常工作。在下面的小节中根据示例进行讨论。\par

\mySmallsection{受限的不变量}

不把C++标准库的全部保证给已移动状态的对象也可以。也就是说，可能有意地限制已移动对象的操作。例如，当对象的有效状态总是需要诸如内存之类的资源时，只有一个部分支持的状态可能会更好地降低移动操作的成本。\par

理想情况下，不支持所有操作的状态是可查的。对象应该知道这个状态，并提供成员函数来检查这个状态。已移动对象也可能拒绝执行此状态下不支持的操作。然而，相应的检查可能会影响性能。\par

C++标准库中，有些类型提供了API来检查对象是否处于移出状态。例如，std::\allowbreak{}futures有一个成员函数valid()，该函数对已移动对象返回false。所以，不同对象检查状态(是否移动)的接口不同。\par

通常，已移动状态是默认构造状态，这意味着已移动状态是不变量的一部分。任何情况下，需要让用户知道什么是可用的，什么不可用。\par


\section{可销毁和可转移}

我们讨论一下如何确保已移动对象满足支持分配和销毁的基本要求。\par

\mySmallsection{6.2.1 可分配和可销毁的已移动对象}

大多数类中，生成的特殊移动成员函数对已移动的对象可以使用赋值操作符和析构函数。如果每个已移动成员都是可赋值和可销毁的，那么整个已移动对象的赋值和销毁都没问题:\par

\begin{itemize}
  \item 通过从相应的源成员中获取状态，赋值将覆盖该成员的未定义状态。
  \item 析构函数将销毁成员(具有未定义状态)。
\end{itemize}

通常应该从对象中移动的成员，但对象处于有效但未定义的状态，所以生成的移动操作通常只对创建成员状态正确的对象进行操作。\par

例如，以下类:\par

\begin{codebox}[cpp]
class Customer {
private:
  std::string name;
  std::vector<int> values;
  ...
};
\end{codebox}

当\textit{customer}移走时，保证\textit{name}和\textit{values}都处于有效的状态，以便析构函数(由\textit{customer}的析构函数调用)可以正常工作:\par

\begin{codebox}[cpp]
void foo()
{
  Customer c{"Michael Spencer"};
  ...
  process(std::move(c));
  // both name and values have valid but unspecified states
  ...
} // destructor of c will clean up name and values (whatever their state is)
\end{codebox}

另外，给\textit{c}赋一个新值也可以，所以就同时赋了\textit{name}和\textit{values}。\par

\mySmallsection{6.2.2 不可销毁的已移动对象}

这时在极少数情况下可能会出现的问题，通常实现赋值操作符或析构函数来完成对成员赋值或销毁的工作。\par

考虑以下类，使用固定大小的线程对象数组，其数量可变，为此在析构函数中显式调用join()来等待线程对象的结束:\par

\begin{codeboxt}[cpp]{basics/tasks.hpp}
#include <array>
#include <thread>
class Tasks {
  private:
  std::array<std::thread,10> threads; // array of threads for up to 10 tasks
  int numThreads{0}; // current number of threads/tasks
  public:
  Tasks() = default;
  // pass a new thread:
  template <typename T>
  void start(T op) {
    threads[numThreads] = std::thread{std::move(op)};
    ++numThreads;
  }
  ...
  // at the end wait for all started threads:
  ~Tasks() {
    for (int i = 0; i < numThreads; ++i) {
      threads[i].join();
    }
  }
};
\end{codeboxt}

目前为止，该类还不支持移动语义，因为用户声明了析构函数。因为禁用复制std::thread，所以也不能复制\textit{Tasks}对象。但是，可以启动多个任务并等待结束:\par

\begin{codeboxt}[cpp]{basics/tasks.cpp}
#include "tasks.hpp"
#include <iostream>
#include <chrono>
int main()
{
  Tasks ts;
  ts.start([]{
    std::this_thread::sleep_for(std::chrono::seconds{2});
    std::cout << "\nt1 done" << std::endl;
  });
  ts.start([]{
    std::cout << "\nt2 done" << std::endl;
  });
}
\end{codeboxt}

现在通过生成默认的移动操作来启用移动语义
(详见\texttt{basics/\allowbreak{}tasksbug.cpp}):\par

\begin{codebox}[cpp]
class Tasks {
private:
  std::array<std::thread,10> threads; // array of threads for up to 10 tasks
  int numThreads{0}; // current number of threads/tasks
public:
  ...
  // OOPS: enable default move semantics:
  Tasks(Tasks&&) = default;
  Tasks& operator=(Tasks&&) = default;
  // at the end wait for all started threads:
  ~Tasks() {
    for (int i = 0; i < numThreads; ++i) {
      threads[i].join();
    }
  }
};
\end{codebox}

这种情况下会遇到麻烦，因为默认生成的移动操作可能会创建无效的任务。比如:\par

\begin{codeboxt}[cpp]{basics/tasksbug.cpp}
#include "tasksbug.hpp"
#include <iostream>
#include <chrono>
#include <exception>
int main()
{
  try {
    Tasks ts;
    ts.start([]{
      std::this_thread::sleep_for(std::chrono::seconds{2});
      std::cout << "\nt1 done" << std::endl;
    });
    ts.start([]{
      std::cout << "\nt2 done" << std::endl;
    });
    // OOPS: move tasks:
    Tasks other{std::move(ts)};
  }
  catch (const std::exception& e) {
    std::cerr << "EXCEPTION: " << e.what() << std::endl;
  }
}
\end{codeboxt}

开始时，通过将两个任务传递给\textit{Tasks}对象\textit{ts}来启动它们。因此，在\textit{ts}中线程数组有两个元素，\textit{numThreads}为2。但是，容器的移动操作会移动元素，因此在\textit{std::move()}之后，\textit{ts}不再包含任何运行的线程对象。因此，\textit{numThreads}只是复制而已，也就是创建了一个无效的任务。析构函数最终会对前两个任务调用join()，这会抛出一个异常(这将导致析构函数中出现致命错误)。\par

目前的问题是，需要两个成员一起才能定义一个有效的状态，而默认生成的移动语义会产生不一致的状态，因此会让析构函数失败。但不可能总是避免这个问题，因为对象销毁时，可能需要显式地对象元素的子集做一些操作。这样，默认生成的移动操作可能不起作用，应该禁用或修复它们。\par

可能的使用办法是:\par

\begin{itemize}
  \item 使用析构函数处理已移动状态(本例中，如果线程对象的joinable()为true，就只能调用join())
  \item 自己实现移动操作
  \item 禁用移动语义(这是特殊移动成员函数的行为)
\end{itemize}

根据规则0，应该将容易出错的资源管理封装在helper类型中，以便应用程序开发者实现0号特殊成员函数。这样，就可以使用一个助手类(模板)，其既提供了成员(std::array和用于实际使用的元素数量的成员)，又提供了移动操作的正确实现。\par

整个问题也是std::thread类设计错误的副作用，该类型不遵循RAII原则。对于所有运行线程的std::thread，必须在调用析构函数之前调用join()(或detach());否则，线程的析构函数将抛出异常。C++20中，可以使用std::jthread类，如果对象为一个正在运行的线程，会自动调用join()。\par


\section{处理已破坏的不变量}

已移动对象可能会破坏“有效但未定义的状态”，这比破坏可销毁要容易得多。我们可以将对象带入一种状态，从而破坏它们的不变量。\par

幸运的是，只有在显式请求移动语义时才会出现这个问题，因为临时对象无论如何都会立即销毁。然而，明确的移动请求不仅仅是用\textit{std::move()}标记对象。可以用以下方法创建已移动对象:\par

\begin{itemize}
  \item 对对象使用\textit{std::move()}
  \item 移动算法(\textit{std::move()}和\textit{std::\allowbreak{}move\_\allowbreak{}backward(\allowbreak{})\allowbreak{}})
  \item 通过将“未删除”的元素移到前面来“删除”元素的算法(例如:\textit{std::\allowbreak{}remove(\allowbreak{})\allowbreak{}}，\textit{std::\allowbreak{}remove\_\allowbreak{}if(\allowbreak{})\allowbreak{}}，\textit{std::\allowbreak{}unique(\allowbreak{})\allowbreak{}})
  \item 移动迭代器
\end{itemize}

如果类的不变量被(生成的)移动操作破坏，有以下解决方法:\par

\begin{itemize}
  \item 修复移动操作，使已移动对象进入不变量稳定的状态。
  \item 禁用移动语义。
  \item 将定义已移动状态的不变量放宽为有效。特别是，成员函数和使用对象的函数必须以不同的方式实现，以处理新的可能状态。
  \item 提供一个成员函数来检查“不变量”的状态，这样类型的用户在用\textit{std::move()}标记后就不会使用这种类型的对象了(或者只使用一组有限制的操作)。
\end{itemize}

来看一些不变量破坏的例子，并讨论如何修复它们。\par

\mySmallsection{6.3.1 破坏由移动值破坏的不变量}

移动操作破坏不变量的第一个原因，与成员的已移动状态就是一个问题，因为该状态是有效的，但不应该发生。\par

考虑纸牌游戏中的Card类。假设每个对象都是一张有效的牌，比如红桃8或方块K。还假设由于某种原因，该值是字符串，并且该类的不变式是每个对象都有一个表示有效卡片的状态。这意味着可能没有默认构造函数，而初始化构造函数会断言该值是有效的。例如:\par

\begin{codebox}[cpp]
class Card {
private:
  std::string value; // rank + "-of-" + suit
public:
  Card(const std::string& v)
  : value{v} {
    assertValidCard(value); // ensure the value is always valid
  }
  std::string getValue() const {
    return value;
  }
};
\end{codebox}

这个类中，生成的特殊移动成员函数创建了一个无效状态，它破坏了类的不变条件，即花色和数值(比如“红心皇后”)。\par

只要我们不使用\textit{std::move()}或其他移动操作，这就不是问题(该类型的析构函数可以从字符串中移动)，但是当调用\textit{std::move()}时，就会遇到麻烦。给新值赋值没什么问题:\par

\begin{codebox}[cpp]
std::vector<Card> deck;
... // initialize deck
Card c{std::move(deck[0])}; // deck[0] has invalid state
deck[0] = Card{"ace-of-hearts"}; // deck[0] is valid again
\end{codebox}

然而，打印已移动卡片的值可能会失败:\par

\begin{codebox}[cpp]
std::vector<Card> deck;
... // initialize deck
Card c{std::move(deck[0])}; // deck[0] has invalid state
print(deck[0]); // passing an object with broken invariant
\end{codebox}

如果print函数的不变量没有破坏，我们会得到一个core dump:\par

\begin{codebox}[cpp]
void print(const Card& c) {
  std::string val{c.getValue()};
  auto pos = val.find("-of-"); // find position of substring (no check)
  std::cout << val.substr(0, pos) << ' '
  << val.substr(pos+4) << '\n'; // OOPS: possible core dump
}
\end{codebox}

此代码可能在运行时失败，因为对于从移动的卡片，不再保证该值包含“-of-”。在这种情况下，\textit{find()}用std::string::\allowbreak{}npos初始化pos，当pos+4用作substr()的第一个参数时，会抛出类型为std::out\_of\_range的异常。\par

完整示例请参见basics/card.hpp和basics/card.cpp。\par

修复方式如下:\par

\begin{itemize}
  \item 禁用移动语义:\par
\begin{codebox}[cpp]
  class Card {
    ...
    Card(const Card&) = default; // disable move semantics
    Card& operator=(const Card&) = default; // disable move semantics
  };
\end{codebox}
  这让移动操作(例如，由std::sort()调用)的开销更大。
  \item 完全禁用复制和移动:\par
\begin{codebox}[cpp]
  class Card {
    ...
    Card(const Card&) = delete; // disable copy/move semantics
    Card& operator=(const Card&) = delete; // disable copy/move semantics
  };
\end{codebox}
  然而，不能再洗牌或对卡片进行排序。
  \item 修复损坏的特殊移动成员函数。\par
  然而，什么是有效的修复(总是分配一个“默认值”，如“梅花A”)？如何确保具有默认值的对象，在不分配内存的情况下性能良好?
  \item 内部允许新状态，但不允许调用getValue()或其他成员函数。\par
  可以对此进行记录(“对于已移动对象，只允许赋值。所有其他成员函数都以对象不处于已移动状态为前提”)，甚至在成员函数内部检查这一点并触发断言或异常。
  \item 纸牌的方式可能没有新状态来扩展不变量。\par
  这意味着必须实现移动特殊成员函数，必须确保对于一个已移动对象，成员值处于这种状态。\par
  通常，已移动状态等同于默认构造状态。因此，这也是提供默认构造函数的机会。理想情况下，还可以提供成员函数来检查这个状态。\par
  这个变化中，这个类的用户必须考虑到字符串的值可能是空，并相应地更新代码。例如:\par
\begin{codebox}[cpp]
  void print(const Card& c) {
    std::string val{c.getValue()};
    auto pos = val.find("-of-"); // find position of substring
    If (pos != std::string::npos) { // check whether it exists
      std::cout << val.substr(0, pos) << ' '
      << val.substr(pos+4) << '\n';
    }
    else {
      std::cout << "no value\n";
    }
  }
\end{codebox}
\end{itemize}

另一方面，getValue()可能返回一个std::optional(自C++17支持)。似乎没有明显的完美解决方案。必须考虑这些修复对程序更大的不变量意味着什么。比如只有一张牌，或者所有的牌都有效等)，然后决定用哪一张。\par

这个类在C++11之前运行良好，C++11之前的标准不支持移动语义(这可能意味着第一个选项是最好的)。因此，C++11可能会为实现类时不可能的类引入状态。这是一种罕见的情况，但移动语义的引入，确实会破坏现有的代码。\par

请参阅Email类的另一个例子，在这个类中，内部标记了已移动状态，以便单独处理，并通过“移除”算法使元素处于已移动状态后，使该状态可见。\par

\mySmallsection{6.3.2 移动的一致性会破坏成员的不变量}

移动操作破坏不变量的第二个原因，与两个成员必须一致但可能被移动破坏的对象有关。正如在线程数组的例子中所看到的，这可能导致破坏析构函数的不一致。通常，析构函数没问题，但自移状态破坏了一个不变量。在类中，有两个不同的值表示形式，一个整数和一个字符串:\par

\begin{codeboxt}[cpp]{basics/intstring.hpp}
#include <iostream>
#include <string>
class IntString
{
  private:
  int val; // value
  std::string sval; // cached string representation of the value
  public:
  IntString(int i = 0)
  : val{i}, sval{std::to_string(i)} {
  }
  void setValue(int i) {
    val = i;
    sval = std::to_string(i);
  }
  ...
  void dump() const {
    std::cout << " [" << val << "/'" << sval << "']\n";
  }
};
\end{codeboxt}

这个类中，确保成员\textit{val}和成员\textit{sval}只是相同值的两种不同表示。这意味着在该类的实现和使用中，通常期望其状态的int和string表示一致。如果在这里使用移动操作，将保留值\textit{val}，但是\textit{sval}不再保证可用\textit{val}的字符串表示。\par

请看下面程序代码:\par

\begin{codeboxt}[cpp]{basics/intstring.cpp}
#include "intstring.hpp"
#include <iostream>

int main()
{
  IntString is1{42};
  IntString is2;
  std::cout << "is1 and is2 before move:\n";
  is1.dump();
  is2.dump();

  is2 = std::move(is1);

  std::cout << "is1 and is2 after move:\n";
  is1.dump();
  is2.dump();
}
\end{codeboxt}

程序有以下输出(已移动符串可能会变成空):\par

\begin{myOutput}
is1 and is2 before move: \\
\ [42/'42'] \\
\ [0/'0'] \\
is1 and is2 after move: \\
\ [42/''] \\
\ [42/'42']
\end{myOutput}

也就是说，自动生成的移动操作打破了两个成员匹配的不变量。\par

这个问题有多大?这至少是个陷阱。同样，您可能会认为在移动之后不应该再使用该值(直到再次设置)。然而，开发者希望使用针对C++标准库对象的策略，该策略声明对象处于有效但未定义的状态。\par

事实是，有了相应的getter，类不再保证int和string的匹配，这可能是类的不变量(隐式或显式声明)。您可能认为最坏的结果是值(现在是未指定的)看起来不同，这取决于如何使用，但使用它没有问题，因为它仍然是有效的int或字符串。\par

然而，依赖于此不变量的代码可能会破坏。该代码可能假设字符串表示至少有一个数字。例如，如果搜索第一个或最后一个数字，肯定会找到。已移动的字符串(通常是空的)不再是这样了。因此，代码不重复检查字符串值中是否有任何字符，可能会遇到未定义行为。\par

同样，如何处理这个问题取决于类的设计者。但是，如果遵循C++标准库的规则，应该使已移动对象处于有效状态，这可能要表示为“没有任何值”的状态。\par

通常，当一个对象的状态具有以某种方式相互依赖时，必须显式地确保已移动状态属于有效状态。但有些情况会出现例外:\par

\begin{itemize}
  \item 对相同值有不同的表示，但是其中一些已移走。
  \item 像\textit{counter}这样的成员对应于成员中的元素数量。
  \item 用布尔值声明字符串的值已验证，但是该验证值已移走。
  \item 所有元素的平均值的缓存值仍然存在，但是值(在容器成员中)已移走。
\end{itemize}

注意，这个类在C++11之前没问题，因为不支持移动语义。当切换到C++11或更高版本，并使用已移动对象时，不变量就破坏了。\par

\mySmallsection{6.3.3 移动的类指针成员破坏不变量}

移动操作破坏不变量的第三个原因，与具有类似指针语义(如(智能)指针)成员的对象有关。\par

考虑以下类的例子，其中对象使用std::shared\_ptr<>共享整型值:\par

\begin{codebox}[cpp]
class SharedInt {
private:
  std::shared_ptr<int> sp;
public:
  explicit SharedInt(int val)
  : sp{std::make_shared<int>(val)} {
  }
  std::string asString() const {
    return std::to_string(*sp); // OOPS: assume there is always an int value
  }
};
\end{codebox}

这个类的对象与副本共享的初始整数值。只要新对象是复制的，一切都好:\par

\begin{codebox}[cpp]
SharedInt si1{42};
SharedInt si2{si1}; // si1 and si2 share the value 42

std::cout << si1.asString() << '\n'; // OK
\end{codebox}

由于只是复制，\textit{SharedInt}成员\textit{sp}总是为它的值分配内存(从std::make\_shared<>()或从已分配内存的现有共享指针复制)。\par

然而，当使用移动语义时，如果使用已移动对象，就会遇到未定义行为:\par

\begin{codebox}[cpp]
SharedInt si1{42};
SharedInt si3{std::move(si1)}; // OOPS: moves away the allocated memory in si1

std::cout << si1.asString() << '\n'; // undefined behavior (probably core dump)
\end{codebox}

问题是在类内部，没有正确处理值可能已移走，因为默认生成的移动操作调用了共享指针的移动操作，该操作将所有权从原始对象移走。这意味着从\textit{SharedInt}的状态移动到成员\textit{sp}不再拥有对象的情况，其成员函数\textit{asString()}无法正确处理。\par

您可能会争辩说，为具有已移动状态的对象调用\textit{asString()}没有意义，因为使用的是未定义的值，但至少标准库保证已移动类型的对象处于有效状态，因此可以调用没有限制的所有操作。如果在用户定义的类型中不提供相同的保证，该类型的用户可能会感到非常惊讶。\par

从鲁棒编程(避免意外、陷阱和未定义行为)的角度来看，我建议您遵循C++标准库的规则。也就是说:移动操作不应该将对象带入破坏不变量的状态。\par

本例中，必须执行以下操作之一:\par

\begin{itemize}
  \item 通过正确处理所有可能的已移动状态来修复类的所有错误操作
  \item 禁用移动语义，在复制对象时没有优化
  \item 显式实现移动操作
  \item 调整并记录类或特定操作的不变量(约束/先决条件)(例如“为已移动对象调用\textit{asString()}是未定义的行为”)。
\end{itemize}

分配内存非常昂贵，最好的解决办法是正确处理所有权移走的情况。这将使用默认构造函数创建对象应具有的状态，这样就可以引入这个状态了。\par

下面的小节将演示这一点和其他修缮的代码。\par

\mySmallsection{修复损坏的成员函数}

第一个选项，修复所有损坏的操作，扩展类的不变量(所有对象的可能状态)，以便所有操作都可以处理已移动状态。我们仍然需要做出设计决策，例如：为已移动对象(或更一般地说，共享指针不拥有整型值的对象)调用asString()时，可以这样:\par

\begin{itemize}
  \item 返回备选值:\par
\begin{codebox}[cpp]
  class SharedInt {
    ...
    std::string asString() const {
      return sp ? std::to_string(*sp) : "";
    }
    ...
  };
\end{codebox}
  \item 抛出异常:\par
\begin{codebox}[cpp]
  class SharedInt {
    ...
    std::string asString() const {
      if (!sp) throw ...
      return std::to_string(*sp);
    }
    ...
  };
\end{codebox}
  \item 在调试模式下强制出现运行时错误:\par
\begin{codebox}[cpp]
  class SharedInt {
    ...
    std::string asString() const {
      assert(sp);
      return std::to_string(*sp);
    }
    ...
  };
\end{codebox}
\end{itemize}

\mySmallsection{禁用移动语义}

第二个选项是禁用移动语义，只使用复制语义。我们前面描述了如何禁用移动语义，必须声明另一个特殊成员函数。通常，对复制特殊成员函数使用\textit{=default}:\par

\begin{codebox}[cpp]
class SharedInt {
  ...
  SharedInt(const SharedInt&) = default; // disable move semantics
  SharedInt& operator=(const SharedInt&) = default; // disable move semantics
  ...
};
\end{codebox}

\mySmallsection{实现移动语义}

第三种选择是实现移动操作，这样就不会破坏类的不变量。\par

为此，必须确定已移动对象的状态应该是什么。为了支持\textit{asString()}可以调用operator*(解引用)而不检查值是否存在，所以必须提供一个值。例如，可以有一个静态的已移动的值，把它赋给那些已移动的对象:\par

\begin{codeboxt}[cpp]{basics/sharedint.hpp}
#include <memory>
#include <string>
class SharedInt {
  private:
  std::shared_ptr<int> sp;
  // special “value” for moved-from objects:
  inline static std::shared_ptr<int> movedFromValue{std::make_shared<int>(0)};
  public:
  explicit SharedInt(int val)
  : sp{std::make_shared<int>(val)} {
  }
  std::string asString() const {
    return std::to_string(*sp); // OOPS: unconditional deref
  }
  // fix moving special member functions:
  SharedInt (SharedInt&& si)
  : sp{std::move(si.sp)} {
    si.sp = movedFromValue;
  }
  SharedInt& operator= (SharedInt&& si) noexcept {
    if (this != &si) {
      sp = std::move(si.sp);
      si.sp = movedFromValue;
    }
    return *this;
  }
  // enable copying (deleted with user-declared move operations):
  SharedInt (const SharedInt&) = default;
  SharedInt& operator= (const SharedInt&) = default;
};
\end{codeboxt}

注意，我们必须对复制特殊成员的函数使用\textit{=default}，因为当我们有声明特殊成员移动函数时，这些函数的默认版本会删除。\par


\section{总结}

\begin{itemize}
  \item 每个类需要说明对象的移动状态，必须确保是可销毁的(通常情况下没有实现特殊成员函数)。然而，类的用户可能期望/要求更多。
  \item C++标准库函数的要求也适用于已移动对象。
  \item 生成的特殊移动成员函数可能会将已移动对象带入某种状态，从而破坏类的不变量。这种情况可能会发生，特别是当:
  \begin{itemize}
    \item[-] 类没有具有确定值的默认构造函数(因此没有自然的移出状态)
    \item[-] 成员的值有限制(比如断言)
    \item[-] 成员的值相互依赖
    \item[-] 使用了具有类似指针语义的成员(指针、智能指针等)。
  \end{itemize}
  \item 如果已移动状态中断不变量或使操作无效，应该使用以下选项来修复这个问题:
  \begin{itemize}
    \item[-] 禁用移动语义
    \item[-] 修缮了移动语义的实现
    \item[-] 类内部处理已破坏的不变量，并将对外隐藏
    \item[-] 通过记录已移动对象的约束和前置条件来消除对类不变量的影响
  \end{itemize}
\end{itemize}
